% Source: https://artofproblemsolving.com/community/c3h2026930_rectangle_and_parallel_lines
% Stable identity: aops-topic-2026930-q01
Place $A=(0,0)$, $B=(l,0)$, $C=(l,w)$, $D=(0,w)$; by symmetry we may take $l>w$. Each bisector makes $45^\circ$ with the sides:
\begin{itemize}
\item from $A$: $y=x$; from $B$: $y=l-x$; from $C$: $y-w=x-l$; from $D$: $y-w=-x$.
\end{itemize}
Intersecting,
\[
P=\left(\dfrac {l}{2},\dfrac {l}{2}\right),\quad Q=\left(l-\dfrac {w}{2},\dfrac {w}{2}\right),\quad R=\left(\dfrac {l}{2},w-\dfrac {l}{2}\right),\quad S=\left(\dfrac {w}{2},\dfrac {w}{2}\right).
\]
\begin{asy}
usepackage("amsmath");
size(6cm);
defaultpen(fontsize(10pt));
real l=6, w=4;
pair A=(0,0), B=(l,0), C=(l,w), D=(0,w);
pair P=(l/2,l/2), Q=(l-w/2,w/2), R=(l/2,w-l/2), T=(w/2,w/2);
draw(A--B--C--D--cycle,linewidth(0.9));
draw(A--P,gray); draw(B--P,gray); draw(C--R,gray); draw(D--R,gray);
draw(P--Q--R--T--cycle,blue+linewidth(0.9));
draw(P--R,dashed+red); draw(T--Q,dashed+red);
dot(A);dot(B);dot(C);dot(D);dot(P);dot(Q);dot(R);dot(T);
label("$A$",A,SW); label("$B$",B,SE); label("$C$",C,NE); label("$D$",D,NW);
label("$P$",P,N); label("$Q$",Q,E); label("$R$",R,S); label("$S$",T,W);
label("$l$",(l/2,0),S); label("$w$",(0,w/2),W);
label("\mbox{}",(min(currentpicture,user=true)+max(currentpicture,user=true))/2,black); // keeps the text after the figure black when LaTeX typesets the labels inline
\end{asy}

The diagonal $PR$ is vertical with length $\dfrac {l}{2}-\left(w-\dfrac {l}{2}\right)=l-w$, and the diagonal $SQ$ is horizontal with length $\left(l-\dfrac {w}{2}\right)-\dfrac {w}{2}=l-w$. They have the common midpoint $\left(\dfrac {l}{2},\dfrac {w}{2}\right)$, so $PQRS$ is a square with diagonal $l-w$ and side
\[PQ=\dfrac{l-w}{\sqrt2}.\]
From $PQ=2019$,
\[|l-w|=\boxed{2019\sqrt2}.\]
